

Following \cite{simo1986variational,simo1990class,simo1992geometrically}
consider a three-field variational formulation of elasticity, with strain fields of the form 
\begin{equation*}
\bm{E} = \hat{\bm{E}}(\hat{\bm{x}}) + \tilde{\bm{E}}
\end{equation*}
where \(\hat{\bm{E}} = \hat{\bm{\varepsilon}}\symbun\FrameBasis\) is the compatible strain tensor derived from the displacement field \cref{eq:finite-alignment-field}, and we refer to $\tilde{\bm{E}}$ as the enhanced part of the strain field. 
This may be decomposed as
\begin{equation*}
\tilde{\bm{E}} = \tilde{\bm{\varepsilon}} \symbun \FrameBasis + \PlaneStrainTensor
\end{equation*}
where \(\PlaneStrainTensor\) is the part of the strain tensor with components on \(\PlaneBasis_{\alpha}\otimes \PlaneBasis_{\beta}\), which is solved for by enforcing the stress constraints:
\[
\PlaneBasis_{\alpha} \cdot \StressTensor \PlaneBasis_{\beta} = 0
\]
Our interest is in the enhaced strain vector \(\tilde{\StrainVector}\). 


\iffalse
\subsection{Configuration Manifolds}
In the context of nonlinear solid mechanics, a body B is embedded in an ambient space A by a map \(\chi \in C\), where C is an infinite dimensional manifold of maps.

The configuration manifold for the geometrically exact rod lacks a vectorial structure. 
For prescribed configuration boundary values \((\tilde{\boldsymbol{x}}, \tilde{\FrameRotation})\) at
\(\xi \in \partial \FrameBody\), the configuration manifold has the representation:
\begin{equation*}
\mathscr{C} \triangleq \left\{\,
  (
  \TraceLocation, 
  \FrameRotation,
  \bm{\alpha}
  ) \colon [0, L] \rightarrow \mathbb{R}^3 
  \times \operatorname{SO}(3)
  \times \mathbb{R}^{n_w}
\bigm|
\chi = (\tilde{\boldsymbol{x}}, \tilde{\FrameRotation})
\text{ on } \partial\FrameBody
\,\right\}.
\end{equation*}
The tangent space associated with a configuration \(\chi\) is represented by: %to be
\begin{equation*}
T_\chi \mathscr{C} \triangleq \left\{
(\boldsymbol{\eta}_x, \boldsymbol{\eta}_{\scriptscriptstyle{\Lambda}})\colon [0, L] \rightarrow \mathbb{R}^3 \times T_{\scriptscriptstyle{\Lambda}}^r \mathrm{SO}(3)
\bigm|
(\boldsymbol{\eta}_x, \boldsymbol{\eta}_{\scriptscriptstyle{\Lambda}})=(\mathbf{0}, \mathbf{0})
\text{ on } \partial\FrameBody
\,\right\}
\end{equation*}
where the right representation \(T_{\scriptscriptstyle{\Lambda}}^r \mathrm{SO}(3)\) of the tangent space on \(\mathrm{SO}(3)\) is used \citep{perez2024nonlinear}.
\fi 

% In the kinematic constitutive framework outlined above, the space \(\mathscr{E}\) of compatible linearized strain tensors \(\hat{\StrainTensor}\) is a \(6+n_w\) dimensional space. 
% The state of strain \(\StrainTensor(\varsigma)\) at all points \(\varsigma\in\PlaneManifold\) is determined completely by the \(6+n_w\) strain parameters \(\bm{e} = (\bm{\gamma},\,\bm{\kappa},\,\bm{\alpha},\,\bm{\alpha}')\). 
% Thus, even before introducing an approximating finite element subspace, it is evident that \(\mathscr{E}\) is  a far weaker space than the infinite-dimensional space of strains allowed for by the full 3D theory.
% In the assumed strain framework, we seek to enrich this space.



\subsection{Enhanced Principle}
Let $\operatorname{B}^*\left[\mathcal{E}\right] \subset \mathscr{E}$ be the space of compatible linear strain tensors generated by the space $\mathcal{E}$ of rod strain fields \(\bm{e}\); i.e. 
$$
\begin{aligned}
\operatorname{B^*}\left[\mathcal{E}\right]
% &=\left\{\bm{E}^h=\sum_{e=1}^{n_{\mathrm{clm}}} \chi_e \bm{\varepsilon}_e(\xi) \symbun\FrameBasis
% \; \bigm| 
% \bm{\varepsilon}_e(\xi)=\sum_{I=1}^{n_{\text {node }}^e} \mathbf{u}_I^e \otimes \operatorname{B^*}[\bm{e}^h(\xi)]
% \right\} \\
&\triangleq \left\{
\hat{\StrainVector} \symbun \FrameBasis \in \mathscr{E} 
\bigm| 
%\hat{\StrainVector}^h=\sum_{e=1}^{n_{\text {elm }}} \hat{\StrainVector}_e(\xi) \chi_e, \text { with } 
\hat{\StrainVector}%_e(\xi)
=\TotalStrainMatrix(\xi,\hat{\rho}) \bm{e},
\quad \bm{e} \in \mathcal{E}
\right\}
\end{aligned}
$$


% in the context of nonlinear solid mechanics, a body B is embedded in an ambient space A by a map $\chi \in C$, where $C$ is an infinite dimensional manifold of maps $\{\chi: B \rightarrow A\}$. 
% In the abscence of constraints (eg, shell and rod theories) $C$ may generally be regarded as a (Linear) Hilbert or Sobolev space. 
% The space of spatial strain measures may be constructed: $\{E: \chi(B)->\chi(B)...\}$.

% \begin{equation}
% \Pi(\mathbf{u}, \tilde{\bm{E}}, \bm{T}) 
% :=\int_0\left.W\left(\bm{x}, \bm{E}_x+\tilde{\bm{\varepsilon}}\right)-\bm{T} \cdot \tilde{\bm{E}}
% \right. \,\mathrm{d} V
% -\int_0 \rho_0 \bm{b} \cdot \mathbf{u} \mathrm{~d} V-\int \tilde{\mathbf{t}} \cdot \mathbf{u} \, \mathrm{d} \Gamma
% \end{equation}

Define the finite dimensional space generated by assumed strain interpolations \(\EnhanShapes\bigl(\xi,\hat{\rho}\bigr)\)
%subspace $\tilde{\mathscr{E}}^h \subset \tilde{\mathscr{E}}$ by the expression
% \begin{equation}
% \tilde{\mathscr{E}}^h
% =\left\{\tilde{\bm{H}}^h \in \tilde{H}
% \;: \tilde{\bm{H}}^h=\sum_{e=1}^{n_{e \mathrm{~lm}}} \tilde{\bm{H}}_e(\xi) \chi_e ; 
% \;
% \tilde{\bm{H}}_e(\xi)=\sum_{I=1} \tilde{\alpha}_I \otimes \mathbf{G}_e^I \text { for } \tilde{\alpha}_I \in \mathbb{R}^{n_{\mathrm{dim}}}
% \right\}
% \label{eq:enhanced-space}
% \end{equation}
\begin{equation*}
\tilde{\mathscr{E}}
\triangleq 
\left\{
\tilde{\StrainVector} \symbun \FrameBasis \in \tilde{\mathscr{E}} \mid %\tilde{\StrainVector}^h=\sum_{e=1}^{n_{\text {lim }}} \tilde{\gamma}_e(\xi) \chi_e, 
%\text { with } 
\tilde{\bm{\varepsilon}}(\xi,\hat{\rho})=\EnhanShapes(\xi, \hat{\rho}) \tilde{\bm{\alpha}}, \text { and } \tilde{\bm{\alpha}} \in \mathbb{R}^{\EnhanDim}\right\}
\label{eq:SR16}
\end{equation*}
where $\tilde{\bm{\alpha}} \in \mathbb{R}^{n_a}$ are $\EnhanDim$ local parameters. %, 
% and $\mathbf{G}^I: \PlaneManifold \rightarrow \mathbb{R}^{n_{\mathrm{dim}}}$ are $\EnhanDim$ vectors of prescribed interpolation functions. 
%
% Here, ${\tilde{\bm{\alpha}}}_{\xi}$ are $\EnhanDim$ section strain parameters, and $\EnhanShapes (\PlanePosition)$ is a matrix of $\mathbb{R}^{n_{n+r}} \times \mathbb{R}^{n_\tau}$ prescribed functions with linearly independent columns, which define the enhanced strain interpolation. 
Consider conditions analogous to those of \cite{simo1992geometrically}:
\begin{enumerate}[label=\roman*]
    \item Stability. 
    % The subspaces spanned by the strains $\operatorname{B^*} [N^I]$ and the assumed functions $\mathbf{G}_e^I$ in \cref{eq:SR16} have null intersection. 
%
    The enhanced strain space $\tilde{\mathscr{E}}$ and the compatible strain space \(\operatorname{B^*}\left[\mathcal{E}\right]\) have null intersection
    
    $$
    \tilde{\mathscr{E}} \cap \operatorname{B^*}\left[\mathcal{E}\right]=\varnothing
    $$

    
\item Consistency. % $S^h$ and $\tilde{\mathscr{E}}^h$ are $L_2$-orthogonal. 
    The subspace of admissible stresses, denoted by $\StressSpaceH \subset L$, must be $L_2$ -orthogonal to the space $\tilde{\mathscr{E}}$ of enhanced strain fields. 
\iffalse 
    Restricting our attention to stress fields discontinuous across element boundaries, we define $\StressSpaceH$ as
    
    $$
    \StressSpaceH:=\left\{\bm{T}^h \in L: \bm{T}^h=\sum_{\ell=1}^{n_{\text {elm }}} \bm{T}_{\ell}(\xi) \chi_{\ell}
    \right\}
    $$
    
    where $\bm{T}_{\ell}: \square \rightarrow \mathbb{L}^{n_{\text {dim }}}$ are prescribed frame-invariant interpolation functions, subject to the constant stress conditions introduced below. 
    The $L_2$-orthogonality requirement implies the condition
    
    $$
    \left<\bm{T}, \tilde{\bm{E}}\right>_{\PlaneManifold}
    = \sum_{I=1} \alpha_I^{e^{\mathrm{T}}} \int_{\mathscr{B}_e} \bm{T}_e \mathbf{G}_e^I \mathrm{~d} V=0
    $$
    %%%%%

    Since any $\left(\tilde{\varepsilon}^h, \sigma^h\right) \in \tilde{\mathscr{E}}^h \times S^h$ is discontinuous across element boundaries, (16) and \cref{eq:SR16} imply

    \begin{equation}
    \left<\bm{\sigma} ,\, \tilde{\bm{\varepsilon}} \right>_{\PlaneManifold}
    :=\int_{x_e} \tilde{\bm{\varepsilon}}  \cdot \bm{\sigma} \, \mathrm{d} V
    =\tilde{\bm{\alpha}} \cdot \int_{\PlaneManifold} \EnhanShapes^{\mathrm{t}} \bm{\sigma}  \, \dd A 
    %=\sum_{I=1} \alpha_I^{e^{\mathrm{T}}} \int_{\mathscr{B}_{\ell}} \bm{T}_{\ell} \mathbf{G}_e^I \, \dd A =0
    \label{eq:20}
    \end{equation}
    for any element Gauss point $\ell$.
    Consequently, we have the requirement that
    %
    \begin{equation}
    \int_{\mathscr{B}_e} \bm{T}_e \mathbf{G}_e^I \, \mathrm{d} V=\mathbf{0}, 
    \quad\text { for all } \quad
    \bm{T}^h=\sum_{e=1}^{n_{\mathrm{elm}}} \bm{T}_e(\xi) \chi_e \in \StressSpaceH
    \label{eq:45}
    \end{equation}
    
    As a result of this condition, equation (9) ${ }_2$ is identically satisfied, and the first term in (9) ${ }_3$ vanishes. 
    Hence, Condition~(ii) effectively eliminates the stress field from the finite element equations. 

    for all $\tilde{\alpha}_I \in \mathbb{R}^{n_{\mathrm{dim}}}$ and any $\bm{T}^h \in \StressSpaceH$. 
\fi
    Thus require \(\tilde{\bm{r}}(\bm{\EnhanParam}) = \mathbf{0}\) with
    \begin{equation}
    \tilde{\bm{r}}(\bm{\EnhanParam}) \triangleq \int_{\PlaneManifold \times \FrameBody} \EnhanShapes^{\mathrm{t}} \StressVector \left(\operatorname{B}[\bm{e}] + \EnhanShapes \bm{\EnhanParam} \right) \, \dd A\, \dd \xi
    \label{eq:def-enhan-residual}
    \end{equation}
\item Patch test. The space $\StressSpaceH$ of admissible nominal stress fields must contain piece-wise constant functions. 
In connection with \cref{eq:def-enhan-residual}, this furnishes the condition:
\begin{equation}
\int_{\PlaneManifold \times \FrameBody} \EnhanShapes \, \dd A \, \dd \xi = \mathbf{0}
\label{eq:condiii}
\end{equation}

\end{enumerate}

\iffalse
\paragraph{Remarks}
\begin{itemize}
\item Condition (i) ensures stability of the method for the linearized problem, and precludes interpolations for the enhanced strain field which would result in rank-deficient methods.
    \item 
Given an interpolation space $\tilde{\mathscr{E}}^h$ of enhanced strain fields, the orthogonality Condition~(ii) may produce an overly constrained space $\StressSpace^h$ which would preclude convergence of the method. 
We prevent such a situation by requiring Condition (iii).
\item For the linear theory, Condition~(iii) is closely related to satisfaction of the classical patch test (see \cite{taylor1986patch}).
\end{itemize}
\fi